\section{The full binary effect body}\label{sec:biasedmetric75}
The spectral boundary of a biased measurement has two distinct
endpoints.  A vanishing eigenvalue of one effect need not give a
coherent angular scale: the other outcome can still destroy coherence.
We therefore retain both endpoint coordinates and determine their
joint finite-use geometry.

Let
\begin{equation}\label{eq:biasedbody75}
 \mathcal E=\left\{(a,x)\in\R\times\R^3:
                   0\le a-|x|\le a+|x|\le2\right\},\qquad
 E_{a,x}=\tfrac12(aI+x\cdot\sigma).
\end{equation}
We identify a point $(a,x)$ of $\mathcal E$ with its effect $E_{a,x}$.
Write $\mathcal M_{a,x}$ for the measurement with ordered effects
$E_{a,x},I-E_{a,x}$ and no residual quantum output.  Its maximum and
minimum eigenvalues, spectral gap, and bias are
\[
 p=\tfrac12(a+|x|),\qquad q=\tfrac12(a-|x|),\qquad
 r=p-q=|x|,\qquad b=p+q-1.
\]
Thus $0\le q\le p\le1$ and $|b|+r\le1$.  When $r>0$ put
$u=x/r$, so that $E_{a,x}=pP_u+q(I-P_u)$, where
$P_u=(I+u\cdot\sigma)/2$.  Directions are not identified under
outcome exchange.

For an integer $N\ge1$ set $\tau=N^{-1}$ and define
\begin{align}
 T_N(s,t)&=|s-t|\sqrt{\frac{N}
                   {\max\{s(1-s),t(1-t)\}+\tau}},
                    &&0\le s,t\le1,\label{eq:spectralmodulus75}\\
 V(p,q)&=p(1-p)+q(1-q)=\tfrac12(1-b^2-r^2),\notag\\
 K_N(p,q)&=r\sqrt{\frac{N}{V(p,q)+\tau}}.
                                            \label{eq:biasedweight75}
\end{align}
For two effects with data $(p,q,u)$ and $(p',q',v)$ let
$h\in[0,\pi]$ be the angle between $u,v$ when both spectral gaps
are positive.  If either gap is zero, set the following angular term
equal to zero, without choosing a direction at the repeated eigenvalue:
\begin{equation}\label{eq:biasedcomparison75}
 H_N(E,F)=T_N(p,p')+T_N(q,q')
             +\min\{K_N(p,q),K_N(p',q')\}\,h.
\end{equation}

\begin{theorem}[Finite-use geometry of binary qubit effects]
\label{thm:biasedmetric75}
For every $E,F\in\mathcal E$ and $N\ge1$,
\begin{equation}\label{eq:biasedmetric75}
 \frac1{8192}\min\{1,H_N(E,F)\}
 \le d_N(\mathcal M_E,\mathcal M_F)
 \le\min\{2,3H_N(E,F)\}.
\end{equation}
The distance permits all common reference-assisted adaptive testers
with at most $N$ calls, including feedback and bounded public stopping,
and uses the unhalved trace norm.  The constants are independent of
both spectra, both directions, and the horizon.  In particular,
\begin{equation}\label{eq:biasedoneuse75}
 d_1(\mathcal M_{a,x},\mathcal M_{a',x'})
                   =|a-a'|+|x-x'|.
\end{equation}
The lower comparison is witnessed by pair-dependent nonadaptive
tests.  It does not require a common estimator of all four target
coordinates.
\end{theorem}

The denominator in \eqref{eq:biasedweight75} involves the two
outcome variances.  Replacing it by the distance to the nearest
spectral face would give an incorrect angular scale on a face such
as $q=0$, $0<p<1$.  The proof below supplies both a scalar-overlap
dilation upper bound and explicit finite-distance lower witnesses.

\subsection{The spectral endpoints}
Put
\[
 B_N(s,t)=\|\operatorname{Bern}(s)^{\otimes N}
                    -\operatorname{Bern}(t)^{\otimes N}\|_1.
\]

\begin{lemma}[A Bernoulli estimate at both support endpoints]
\label{lem:bernoulliproduct75}
For $s,t\in[0,1]$,
\begin{equation}\label{eq:bernoulliproduct75}
 \frac1{64}\min\{1,T_N(s,t)\}
       \le B_N(s,t)\le\min\{2,3T_N(s,t)\}.
\end{equation}
For a fixed one of the two parameters, $B_N$ increases as the other
parameter moves away from it on either side.
\end{lemma}
\begin{proof}
Assume $s<t$, put $\Delta=t-s$, and write
$v=\max\{s(1-s),t(1-t)\}$.  The identical pair is immediate.
A coupling of individual bits gives $B_N\le2N\Delta$.
The root fidelity $F=\sqrt{st}+\sqrt{(1-s)(1-t)}$ satisfies
\begin{align*}
 1-F
 &=\frac{\Delta^2}{2}
       \left(\frac1{(\sqrt s+\sqrt t)^2}
            +\frac1{(\sqrt{1-s}+\sqrt{1-t})^2}\right)\\
 &\le\frac{\Delta^2(1+\Delta)}{2t(1-s)}
 \le\frac{\Delta^2}{2v},
\end{align*}
whenever $v>0$.  For the last inequality, both endpoint variances
are at most $t(1-s)/(1+\Delta)$: after cancellation the required
inequalities are $s(1+\Delta)\le t$ and
$(1-t)(1+\Delta)\le1-s$.  They follow respectively from
$\Delta(1-s)\ge0$ and $t\Delta\ge0$.
The product-fidelity inequality gives
$B_N\le2\sqrt N\Delta/\sqrt v$.
For $z=Nv$, the inequality
$\min\{2,2/\sqrt z\}\sqrt{z+1}\le2\sqrt2<3$
combines the two upper bounds.  At $v=0$ use the coupling bound.

For the lower bound, complement both bits if necessary to have
$s+t\le1$.  If $t>3/4$, then $\Delta\ge2t-1>1/2$ and
$B_N\ge B_1=2\Delta>1$, which suffices.  Otherwise $0<t\le3/4$.
Set
\[
 k=\min\left\{N,\max\left\{1,
                           \left\lfloor\frac1{2t}\right\rfloor
                              \right\}\right\},\qquad
 m=\lfloor N/k\rfloor.
\]
Then $k\ge\frac14\min\{N,t^{-1}\}$, $m\ge N/(2k)$, and
$(k-1)t\le1/2$.  The event of no success in a block of $k$ trials
has probabilities $e_s=(1-s)^k$ and $e_t=(1-t)^k$, with
\[
 e_s-e_t\ge k\Delta(1-t)^{k-1}\ge k\Delta/2.
\]
There is a common stochastic processing of this binary event whose
success probability is $d+ce$, where
\[
 c=\frac1{\max\{e_s+e_t,2-e_s-e_t\}}\ge\tfrac12,
 \qquad d=\tfrac12(1-c(e_s+e_t)).
\]
Indeed $d\ge0$, $d+c\le1$, and the processed pair has parameters
$(1\pm c(e_s-e_t))/2$.  Apply
Lemma~\ref{lem:bernoullimajority73} to $m$ independent blocks and
discard the unused calls.  Contraction under this processing gives
\[
 B_N(s,t)\ge\frac1{16\sqrt2}
       \min\{1,\Delta\sqrt{N\min(N,t^{-1})}\}.
\]
Since $v\ge t(1-t)\ge t/4$,
$T_N(s,t)\le2\Delta\sqrt{N\min(N,t^{-1})}$.
This proves the stated lower constant.

For monotonicity, the likelihood ratio of the two binomial count
laws is monotone in the count.  For a fixed larger parameter their
total variation is the maximum, over upper count intervals, of the
larger-parameter probability minus the smaller-parameter probability.
The latter probabilities decrease when the smaller parameter
decreases.  Fixing the smaller parameter gives the analogous
statement with lower count intervals.  Endpoint cases follow by
continuity of the finite sums.
\end{proof}

\begin{lemma}[Spectral projection without cancellation]
\label{lem:spectralprojection75}
For arbitrary directions and spectra,
\begin{equation}\label{eq:spectralprojection75}
 B_N(p,p')\le d_N(\mathcal M_E,\mathcal M_F),\qquad
 B_N(q,q')\le d_N(\mathcal M_E,\mathcal M_F).
\end{equation}
If the two directions agree, then
\begin{equation}\label{eq:alignedspectral75}
 d_N(\mathcal M_{p,q,u},\mathcal M_{p',q',u})
                              \le B_N(p,p')+B_N(q,q').
\end{equation}
Consequently, with $d=d_N(\mathcal M_E,\mathcal M_F)$,
\begin{equation}\label{eq:angularprojection75}
 \begin{split}
 d_N(\mathcal M_{p,q,u},\mathcal M_{p,q,v})&\le3d,\\
 d_N(\mathcal M_{p',q',u},\mathcal M_{p',q',v})&\le3d.
 \end{split}
\end{equation}
\end{lemma}
\begin{proof}
Suppose first $p\ge p'$.  Repeatedly input a maximum-eigenvalue
eigenstate of $E$.  Its success probabilities under the two
hypotheses are $p$ and some $t\le p'$.  The monotonicity in
Lemma~\ref{lem:bernoulliproduct75} gives
$B_N(p,p')\le B_N(p,t)\le d$.  If $p<p'$, interchange the
hypotheses.  For $q\le q'$, use a minimum-eigenvalue eigenstate
of $E$; the other success probability is at least $q'$.  Reverse
the hypotheses if $q>q'$.  These are separate admissible witnesses;
no single input is asserted to witness both endpoint changes.

For aligned directions, first make the common projective measurement
$\{P_u,I-P_u\}$ and, according to its outcome, generate a bit with
success probability $p$ or $q$.  This implements exactly the
measurement channel, including on inputs entangled with a reference:
the reference blocks are the corresponding partial matrix elements
weighted by the displayed success probabilities.  Supply in advance
$N$ independent $p$-bits and $N$ independent $q$-bits; at a call
consume the next unused bit from the selected array.  Every fixed
adaptive tester is a common channel of these two arrays.  The
triangle inequality for tensor products therefore bounds its output
distance by $B_N(p,p')+B_N(q,q')$.  This proves
\eqref{eq:alignedspectral75} with arbitrary references, feedback,
and stopping; all unqueried row bits are discarded. Combining it with
\eqref{eq:spectralprojection75} bounds every aligned spectral
change by $2d$.  The triangle inequality, aligning at $u$ or at $v$,
then proves \eqref{eq:angularprojection75}.
\end{proof}

\subsection{Angular geometry with both eigenvalues retained}
\begin{theorem}[Angular transition for a biased measurement]
\label{thm:biasedangular75}
For fixed $0\le q\le p\le1$, $N\ge1$, and directions at angle
$h\in[0,\pi]$,
\begin{equation}\label{eq:biasedangular75}
 \frac1{512}\min\{1,K_N(p,q)h\}
 \le d_N(\mathcal M_{p,q,u},\mathcal M_{p,q,v})
 \le\min\{2,\sqrt2 K_N(p,q)h\}.
\end{equation}
More precisely, put
\[
 A=\sqrt{p(1-p)}+\sqrt{q(1-q)},\qquad
 z=(p-q)\sin(h/2).
\]
When $A>0$,
\begin{equation}\label{eq:biasedfiniteupper75}
 d_N\le\min\left\{2,2Nz,
       2\sqrt{1-\left(\frac{A^2}{A^2+z^2}\right)^N}\right\}.
\end{equation}
When $A=0$, the bound $\min\{2,2Nz\}$ remains valid.
\end{theorem}
\begin{proof}
The one-use distance is $2z$: the effect difference is traceless
with operator norm $z$, and the same reference trace-norm duality
as in Lemma~\ref{lem:noisyupper73} applies.  Replacing the calls one
at a time gives the adaptive bound $2Nz$.

We next construct a dilation with scalar cross overlap.  Common
unitary conjugation puts the two directions in one real meridian.
Let
\[
 J=\begin{pmatrix}0&-1\\1&0\end{pmatrix},\qquad
 R_t=\exp(tJ),\qquad t=h/2,\qquad
 E=\operatorname{diag}(p,q),\quad F=R_tER_t^{\mathsf T}.
\]
A measurement dilation is
\[
 W_E\psi=|+\rangle_C|+\rangle_L\sqrt E\,\psi
            +|-\rangle_C|-\rangle_L\sqrt{I-E}\,\psi.
\]
The square-root vectors belong to a further two-dimensional
environment.  Tracing both environments leaves the ordered
classical output and implements the measurement even with a
reference.  Put
\[
 f=\sqrt{pq}+\sqrt{(1-p)(1-q)}.
\]
For any $\phi$, apply the common rotation $R_{\phi-t}$ to the
two-dimensional environment of the dilation $W_F$.  The cross
overlap is
\begin{align*}
 W_E^\dagger(I\otimes R_{\phi-t})W_F
 &=\bigl(\sqrt E\,R_\phi\sqrt E
        +\sqrt{I-E}\,R_\phi\sqrt{I-E}\bigr)R_{-t}\\
 &=\bigl(\cos\phi\,I+f\sin\phi\,J\bigr)R_{-t}.
\end{align*}
Here the rotation acts as the identity on the outcome registers.
When $f>0$, choose
\[
 D=\sqrt{f^2\cos^2t+\sin^2t},\qquad
 \cos\phi=\frac{f\cos t}{D},\qquad
 \sin\phi=\frac{\sin t}{D}.
\]
Then the displayed overlap is $cI$, where $c=f/D\in[0,1]$.
This construction also covers $t=0$ and $t=\pi/2$.
To express it in the spectral coordinates, put
$T=\sqrt{p(1-q)}+\sqrt{q(1-p)}$.  Direct multiplication gives
\[
 fT=A,\qquad (1-f^2)T^2=(p-q)^2.
\]
If $A>0$ then $f,T>0$, and therefore
\begin{equation}\label{eq:scalaroverlap75}
 c^2=\frac{f^2}{f^2\cos^2t+\sin^2t}
          =\frac{A^2}{A^2+(p-q)^2\sin^2t}.
\end{equation}

Fix any common adaptive tester.  Purify its initial state and all
intermediate operations, and retain the measurement environments in
this proof.  The environments are never accessed by the tester.
At each call the scalar identity for the cross overlap multiplies
the inner product of the two purifications by $c$, independently of
the current input and reference.  The common intervening isometries
preserve the inner product.  A tester that stops early is padded to
$N$ calls on dummy inputs, with their outcomes ignored in its final
output.  Thus two purifications of that final output have inner
product $c^N$.  Trace-norm contraction and the pure-state formula
give $d_N\le2\sqrt{1-c^{2N}}$, proving
\eqref{eq:biasedfiniteupper75}.  This is a pairwise dilation
comparison; it does not require a processor for the entire effect
body or an unknown-input environment-seizing procedure.

Since $V\le A^2\le2V$ and
$1-(1+y)^{-N}\le Ny$ for $y\ge0$, the two nontrivial upper
bounds imply
\[
 d_N\le\min\{Nrh,\sqrt N\,rh/\sqrt V\}
            \le\sqrt2 K_N(p,q)h
\]
when $V>0$.  The final inequality follows by putting $z=NV$
and using $\min\{1,z^{-1/2}\}\sqrt{z+1}\le\sqrt2$.
If $V=0$, the only possibilities are scalar deterministic effects
or a projective measurement, and the hybrid bound suffices.

For the lower bound, two explicit witnesses suffice.  First suppose
$r>0$ and $h>0$.  Input a pure state in the direction $u-v$.
Its success probabilities are
\[
 \frac{1+b+r\sin(h/2)}2,\qquad
 \frac{1+b-r\sin(h/2)}2.
\]
Their difference is $r\sin(h/2)$.  The larger of their Bernoulli
variances is at most $1-b^2$: after complementing both output laws
we may assume $a=1+b\le1$; both success probabilities lie in
$[0,a]$, so their variances are at most $a\le a(2-a)=1-b^2$.
Lemma~\ref{lem:bernoulliproduct75}, together with
$\sin(h/2)\ge h/\pi$, now gives
\begin{equation}\label{eq:biasedproductlower75}
 d_N\ge\frac1{64\pi}\min\left\{1,
        rh\sqrt{\frac{N}{1-b^2+\tau}}\right\}.
\end{equation}
The repeated pure input is a nonadaptive test.

For the second witness, choose a half-turn of the input qubit about
an axis normal to the plane of $u,v$.  It sends both Bloch
directions to their negatives.  With probability $1/2$ use the
original measurement, and with probability $1/2$ apply this
half-turn before the call and swap the two recorded outcomes
afterwards.  Discard the random choice.  The resulting measurement
has effects $(I\pm r u\cdot\sigma)/2$ or
$(I\pm r v\cdot\sigma)/2$ under the respective hypotheses:
the bias has been removed and the visibility is $r$.
The same random preprocessing and postprocessing are used under
both hypotheses and on every call.  Any test for this unbiased
pair is consequently an admissible test for the original pair.
Theorem~\ref{thm:noisycoding73} gives
\begin{equation}\label{eq:biasedentangledlower75}
 d_N\ge\frac1{64}\min\left\{1,
        rh\sqrt{\frac{N}{1-r^2+\tau}}\right\}.
\end{equation}
The witness is the finite entangled-block test of
Lemma~\ref{lem:noisylower73}, with this common processing.
If the directions coincide or are antipodal, the required
half-turn is chosen about any perpendicular axis.

Finally, feasibility $|b|+r\le1$ implies
\[
 \min\{1-b^2,1-r^2\}\le2(1-b^2-r^2)=4V.
\]
For example, write $x=\max\{|b|,r\}$ and
$y=\min\{|b|,r\}$; then $x\le1-y$, $0\le y\le1/2$, and
$x^2+2y^2\le(1-y)^2+2y^2\le1$, which is the asserted
inequality.  Hence the larger of the two coefficients in
\eqref{eq:biasedproductlower75} and
\eqref{eq:biasedentangledlower75} is at least $K_N(p,q)/2$.
Their stronger lower bound is at least
$(128\pi)^{-1}\min\{1,K_N(p,q)h\}$, which implies
\eqref{eq:biasedangular75}.  If $r=0$ or $h=0$, the two channels
coincide and the assertion is immediate.
\end{proof}

\begin{proof}[Proof of Theorem~\ref{thm:biasedmetric75}]
For positive gaps use a path whose angular leg has the spectrum
with the smaller angular coefficient.  Lemma
\ref{lem:spectralprojection75}, followed by
Lemma~\ref{lem:bernoulliproduct75} and
Theorem~\ref{thm:biasedangular75}, gives
\[
 d_N(\mathcal M_E,\mathcal M_F)
 \le3T_N(p,p')+3T_N(q,q')
             +\sqrt2\min\{K_N(p,q),K_N(p',q')\}h
 \le3H_N(E,F).
\]
If one gap vanishes, that scalar effect may be aligned with the
other direction at no cost, so the same estimate holds with angular
term zero.

For the lower bound, put $A=T_N(p,p')$, $B=T_N(q,q')$, and
$C=\min\{K_N(p,q),K_N(p',q')\}h$.
The spectral projection lemma gives
$d_N\ge64^{-1}\min\{1,A\}$ and
$d_N\ge64^{-1}\min\{1,B\}$.
Its angular projections and Theorem~\ref{thm:biasedangular75}
give $d_N\ge1536^{-1}\min\{1,C\}$.
The largest of these three bounds is at least
$4608^{-1}\min\{1,A+B+C\}$, which proves the weaker displayed
constant $1/8192$.  At a scalar effect only the two spectral bounds
are needed.  To check the assertion about lower witnesses, let
$d_N^{\rm na}$ be the supremum over nonadaptive tests.  The two
spectral lower bounds hold with $d_N^{\rm na}$, and the aligned
spectral upper bound is consequently at most $2d_N^{\rm na}$.
Apply the triangle inequality to the output of a nonadaptive
angular witness on the three-channel path.  Its actual-pair leg
is at most $d_N^{\rm na}$, and its aligned spectral leg is at most
$2d_N^{\rm na}$.  The same lower comparison thus holds with
$d_N^{\rm na}$ in place of $d_N$.

It remains to verify the one-use identity.  For a reference-assisted
input state $\rho$, the difference has output blocks $B_\rho$ and
$-B_\rho$, where
$B_\rho=\tr_A[((E-F)\otimes I)\rho]$.  Duality against
Hermitian contractions gives
$\|B_\rho\|_1\le\|E-F\|_{\rm op}$.  An extremal eigenstate
of $E-F$ attains equality.  Thus
$d_1=2\|E-F\|_{\rm op}$; the two eigenvalues of $E-F$ are
$((a-a')\pm|x-x'|)/2$.  This proves
\eqref{eq:biasedoneuse75}.
\end{proof}

\begin{remark}[Boundary faces and repeated eigenvalues]
\label{rem:biasedfaces75}
On $p=q=t$, the measurement is the input-independent coin
$\operatorname{Bern}(t)$ and all directions represent the same
effect.  Its exact finite-use distance along this scalar interval is
$B_N(t,t')$.  On the face $q=0$, $0<p<1$, the angular coefficient is
\[
 K_N(p,0)=p\sqrt{\frac{N}{p(1-p)+N^{-1}}}.
\]
For a fixed $p$ in this interval it has square-root-horizon order,
despite the zero eigenvalue.  At the projective edge $(p,q)=(1,0)$
it equals $N$.  At the two deterministic apices $p=q=0$ and
$p=q=1$, the angular coefficient vanishes while the spectral
moduli retain their finite-use support cutoff.  These cases are
part of the same theorem, without a spectral margin imposed on
the target.
\end{remark}
